\documentclass[12pt]{article}

\usepackage[a4paper, total={6in, 8in}]{geometry}
\usepackage{textcomp}

\begin{document}
\setlength{\parindent}{0pt}

\def \t {\textrightarrow}
\def \v {\vspace{1em}}
\def \bi {\begin{itemize}}
\def \ei {\end{itemize}}
\def \s[#1] {\section*{#1}}
\def \ss[#1] {\subsection*{#1}}
\def \sss[#1] {\subsubsection*{#1}}

\s[Two war poems]
Both author fought during the war, and reflected on their experience at the front in some of their poems

\ss[The soldier - Rupert Brooke]
He was part of a poetical movement known as "georgeans" \t he is a "georgean poet" \t they wrote their poems between 1910 and 1915, when George V was king of the UK \\
The georgeans tended to write about the english countryside \t they considered it as the essence of "englishennes" \t influence of Wordsworth \\
Robert Brooke wrote this before leaving for the front \t he is enthusiastic and is a volounteer \t he died some weeks later after writing this

\v

This poem is an italian sonnet \t made of 14 lines, grouped into an octave and a sestet (the english sonnet is made of 3 quatrains and a couplet) \\
At the beginning of the 20th century started to write poems in free verses (no forms, no rhymes) \t howerver georgean poems were traditional

\v

First line \t assonance in "o" and allitteration in "d" \\
Second verse in "f", in line 3 there is a cesura \\
Richer dust because it is where an english soldier died \\
Personification of dust \\
Reference to the english countryside \t english / england = is a polyptoton (repetition of words that share the same root) \\
The sestet is more metaphysical \\
Eternal life is a sort of after life \\
This is an idealistic poem

\ss[Dulce et Decorum est - Wilfred Owen]
The title is a partial quotation of the latin poem of Horace (orazio) \t however the quotation is ironic, there is nothing good about war is the message \\
There are rhymes (so no free verses) \t but the verses are iambic pentameters (or maybe soldier in this one not? ask) \\
4 stanzas

\v

First one sets the scene \\
"Like old beggars" is a simile \t he talks about a group of soldiers \t unheroic representation \\
"Like hags" (old women) \t another simile \\
"Sludge" = mud \\
Verse 3 and 4 \t they have gon over the top, they attacked \t now they return to their trenches \\
"Limped on " = zoppicare \\
Line 6 \t diacope of "all"

\v

In the second one, a gas attack is described \\
One was not quich enough to wear the mask \\
"Under a green a sea" \t simile \t the gas turned the air green \t he sees his comrade dying

\v

Stanza 3 is about the trauma, that seeing his comrade dying has created in the writer \\
Consonance of "s", first line \\

\v

Stanza 4 conveys the message of the whole poem \\
"Like a devils'" \t simile, his fase is like the one of a devil \\
Message is that pro war propaganda should not exist \t it convinces young men to leave for the front

\v

The view of line 1 and the one of "my friend" in line 9 make reference to Jessie Pope, a (singoistic = patriotic) poetess and journalist \\
She became famous during the war for writing several pro war poems and articles on british newspapers, that wanted to convince young british men to go to war \\
This poem had another title: "To a certain poetess" \t then it was changed, but because of this change some scholars believe that "my friend" and "you" can be referred to members of the propaganda system in general

\v

Owen had a similar but different story to Brooke \t he was a volounteer too, but then he developed an antiwar position \\
He got a severe form of shell-shock \t he was hospitalized at Craiglockhart hospital (scottish hospital where shell shcoks were sent) \t after some months of treatment, they sent him to the front \\
He died only one week before the end of the war

\s[1920s - 1930s in England]
These years have 3 labels:
\bi
  \item the interwar years \t purely historical
  \item the age of anxiety \t intellectual atmosphere characterzing these years
  \item anglo-american modernism \t most relevant cultural movement that developed during this period
\ei
The UK and US lived these years in very different ways for ? it was tuff \t british economy started to have economic problems, and during the inter war years the british empire started collapsing \\
The real collapse happened after the WWII, but slowly starts here \t in 1921 the biriths empire lost Ireland, its nearest colony \t Eire, the Repulic of Ireland was established \\
Northern Ireland remains part of the UK \\
Now the so known "irish troubles" (1921 - 1994) begin \t in 1994 a truce was signed (not a peace treaty) \\
The irish troubles consisted in a civil war between northern irish catholics and northern irish anglicans \t catholics wanted northern ireland to become part of the Eire, while anglicans wanted to remain part of the UK \\
Many riots and terroristic attacks \t that happened in the UK too \\
The organization that organized these attacks was the I.R.A \t irish republic army (catholic terroristic org.), that fought for making norther ireland part of Eire \\
In the 1930s, according to the Statute of Westminister, some colonies were given the state of "dominions" \t not independent, but could be more autonomous \t the UK was not able to controll its empire with an iron fist anymore \t they had to make some concessions \\
This does not happen however with India \t after the WII, with the process of decolonization, the english empire got dismantelled 
\end{document}
